% Volume VII: Formal Synthesis
% Part XIII. A Mathematics of Disciplined Distrust
% Chapter 76: The Freedom Region
% Spine: proof-spines/ch076-spine.md

\chapter{The Freedom Region}
\label{ch:ch076}

This chapter defines the \textbf{freedom region} as a feasible institutional set rather than a single score. The region is
\[
\mathfrak F
=
\left\{
(\kappa,\rho,\phi,\theta):
L_{\mathrm{domination}}
+
L_{\mathrm{exposure}}
+
L_{\mathrm{error}}
+
L_{\mathrm{inaction}}
\leq \varepsilon
\right\}.
\]
\ToyModel\ The parameters collect jointly maintained constraints, while the losses register distinct failure modes: domination, exposure, error, and inaction. A society counts as free, on this view, when it remains inside the feasible region rather than at one privileged numerical optimum.

The importance of \(\mathfrak F\) is that it refuses to collapse liberty into a commensurable scalar. Different orders may occupy different points of the region while preserving the same basic condition: persons are not forced into exposure, domination, or reckless institutional response. Freedom is thus modeled as a region of viable joint constraints, completing Part XIII's claim that disciplined distrust exists to protect a structured social space rather than to maximize suspicion or transparency in isolation.
